\documentclass{article}
\usepackage{eason}

\usepackage{pdfpages}

\usepackage[a4paper, margin=1in]{geometry}

\usepackage{lipsum}
\usepackage{mwe}
\usepackage{environ}

% \RenewEnviron{figure}{}
% \RenewEnviron{table}{}
% \RenewEnviron{listing}{}

\usemintedstyle{bw}

\titleformat{\section}{\normalfont\Large\bfseries}{}{0pt}{}

\newcommand{\rowdots}{\vphantom{0}\smash{\raisebox{-0.7ex}{\vdots}}}

\begin{document}
% \maketitle

{\LARGE\bf Project}

\vspace*{-10mm}
\hfill\framebox{\parbox[t][5 true cm][c]{11 true cm}{\hfil
        Space for the project label }}

\bigbreak

\begin{center}\LARGE\bf
  Project Title
\end{center}

\section{Question}

\autoref{eq:euler} is \term{Euler's identity}, and is illustrated in \autoref{fig:example}.

\begin{equation}
    \label{eq:euler}%
    e^{i\pi} + 1 = 0.
\end{equation}

\begin{figure}[ht!]
    \centering
    \includegraphics[width=0.5\textwidth]{example-image-a}
    \caption{Example Image}
    \label{fig:example}
\end{figure}

\section{Task}

The output is in \autoref{tab:example}.The code is as in \autoref{lst:example}.

\begin{table}[ht!]
    \centering
    \caption{Example Table}
    \smallskip 
    \begin{tabular}{|c|c|c|}
        \hline
        Column 1 & Column 2 & Column 3 \\
        \hline
        Data 1 & Data 2 & Data 3 \\
        Data 4 & Data 5 & Data 6 \\
        \hline
    \end{tabular}
    \label{tab:example}
\end{table}

\clearpage
\appendix
\section{Listings}

Program was coded in a \verb|.ipynb| Jupyter Notebook (\verb|notebook.ipynb|, a fresh version \verb|notebook-fresh.ipynb| without outputs) using Python 3.15.0, utilising packages \verb|numpy|, \verb|pandas| and \verb|matplotlib|. In addition, a custom \verb|macros.tex| is used for producing the \LaTeX formulae in titles in the diagrams, but is itself irrelevant to the code and the mathematics.

A program \verb|convert-csv-to-latex.py| is used to convert the exported \verb|.csv| data files to \LaTeX-formatted \verb|tabular| environments, and is itself irrelevant to the mathematical contents.

The following listings are key segments of the original code (exported into \verb|notebook.py|) related to the problem, and might not run individually. The Python code can be run as a whole, and will export the data into \verb|.csv| and figures into \verb|.png|.

All programs mentioned above apart from \verb|macros.tex| and \verb|convert-csv-to-latex.py| are submitted for inspection.

\begin{listing}[htp]
    \inputminted{python}{src/notebook.py}
    \caption{Testing Python code}
    \label{lst:example}%
\end{listing}

\end{document}